\documentclass[11pt]{article}
\usepackage[margin=1in]{geometry}
\usepackage{amsmath,amssymb,amsthm}
\usepackage{graphicx}
\usepackage{microtype}
\usepackage{titlesec}
\usepackage[hidelinks]{hyperref}
\titleformat{\section}{\normalfont\Large\bfseries}{\thesection}{0.8em}{}

\newtheorem{theorem}{Theorem}
\newtheorem{corollary}[theorem]{Corollary}
\newtheorem{remark}[theorem]{Remark}
\newcommand{\R}{\mathbb{R}}
\newcommand{\IS}{\operatorname{IS}}
\newcommand{\as}{\operatorname{as}}
\newcommand{\vrad}{\operatorname{vrad}}
\newcommand{\Cin}{\mathcal{C}_{\mathrm{in}}}

\title{A radial-envelope extension for extremal general affine surface area}
\author{Candidate partial result for arXiv:2103.00294}
\date{19 July 2026}

\begin{document}
\maketitle

\noindent\textbf{Status.} Candidate partial result, likely valid, pending human review.

\section{Source problem and claimed scope}

Hoehner studies extremal general affine surface areas for centered convex
bodies and function classes determined by global monotonicity of
$\varphi(t)/\sqrt t$ \cite{Hoehner}.  At the end of Section 7.4 (page 21), the
paper asks to ``obtain the results of this article for a richer class of
functions than $\operatorname{Conc}^{-}(0,\infty)\cup
\operatorname{Conc}^{+}(0,\infty)$.''

\begin{center}
\includegraphics[width=\textwidth]{figures/open_problem_crop.png}
\end{center}

This note gives a partial answer.  It strictly enlarges the function class for
which the inner maximal functional is finite, attained, and continuous, and it
replaces the source's endpoint radial estimate by a sharp radial envelope.
It does not extend every theorem of the source paper.

\section{Definitions}

Let $B_n$ be the Euclidean unit ball in $\R^n$.  For a convex body $K$,
write
\[
  \vrad(K)=\left(\frac{|K|}{|B_n|}\right)^{1/n}.
\]
All convex bodies below have centroid at the origin.  Let $\Cin(K)$ denote the
convex bodies $C\subset K$ whose centroid is the origin.

Let $\operatorname{Conc}(0,\infty)$ be the source class of positive concave
functions $\varphi:(0,\infty)\to(0,\infty)$ satisfying
$\varphi(0)=0$ and $\varphi(t)/t\to0$ as $t\to\infty$.  Its general affine
surface area is denoted by $\as_\varphi$.  Define
\[
  \IS_\varphi(K)=\sup_{C\in\Cin(K)}\as_\varphi(C).
\]
The two source facts used below are the affine isoperimetric inequality
\begin{equation}\label{eq:isoper}
 \as_\varphi(C)\le
 \as_\varphi(\vrad(C)B_n)
\end{equation}
for centered $C$, and upper semicontinuity of $\as_\varphi$ under Hausdorff
convergence \cite{Ludwig,Hoehner}.  On a Euclidean ball,
\begin{equation}\label{eq:ball}
 \as_\varphi(rB_n)
 =r^n\varphi(r^{-2n})|\partial B_n|
 =n|B_n|\frac{\varphi(t)}{\sqrt t},
 \qquad t=r^{-2n}.
\end{equation}

Put
\[
 h_\varphi(t)=\frac{\varphi(t)}{\sqrt t},\qquad
 L_\varphi=\limsup_{t\to\infty}h_\varphi(t),\qquad
 H_\varphi(a)=\sup_{t\ge a}h_\varphi(t).
\]
Call $\varphi$ \emph{tail-gapped} when
\begin{equation}\label{eq:tailgap}
 \begin{aligned}
 L_\varphi&<\infty,\\
 H_\varphi(a)&>L_\varphi\quad\text{for every }a>0.
 \end{aligned}
\end{equation}

\section{Proof intuition}

The source proof controls an inscribed body $C$ only through its volume
radius.  Formula \eqref{eq:ball} shows that the correct one-dimensional object
is not an endpoint value of $h_\varphi$, but its tail envelope $H_\varphi$.
Finite tail limsup makes this envelope finite.  The strict tail gap then
prevents a maximizing sequence from collapsing: a degenerate sequence can
have limiting affine surface area at most $n|B_n|L_\varphi$, while a sufficiently
small centered ball inside $K$ has strictly larger value.  Once collapse is
excluded, Blaschke selection and upper semicontinuity give an extremizer.
Continuity follows separately from inclusion monotonicity and the elementary
dilation estimate, so it does not require global monotonicity of
$h_\varphi$.

\section{Radial-envelope theorem}

\begin{theorem}\label{thm:main}
Let $\varphi\in\operatorname{Conc}(0,\infty)$ be tail-gapped.  Then for every
centered convex body $K\subset\R^n$:
\begin{enumerate}
 \item $0<\IS_\varphi(K)<\infty$, and some
 $C_{\max}\in\Cin(K)$ satisfies
 $\IS_\varphi(K)=\as_\varphi(C_{\max})$;
 \item the sharp radial-envelope bound holds:
 \begin{equation}\label{eq:envelope}
  \IS_\varphi(K)
  \le n|B_n|\,H_\varphi\bigl(\vrad(K)^{-2n}\bigr)
 \end{equation}
 \item equality in \eqref{eq:envelope} holds when $K=RB_n$ for every $R>0$;
 \item $K\mapsto\IS_\varphi(K)$ is continuous on centered convex bodies in
 the Hausdorff metric.
\end{enumerate}
Moreover, every function in the source class
$\operatorname{Conc}^{-}(0,\infty)$ is tail-gapped.
\end{theorem}

\begin{proof}
Set $R=\vrad(K)$ and $a=R^{-2n}$.  For $C\in\Cin(K)$, volume inclusion gives
$\vrad(C)\le R$.  With $t_C=\vrad(C)^{-2n}\ge a$, equations
\eqref{eq:isoper} and \eqref{eq:ball} give
\[
 \as_\varphi(C)
 \le n|B_n|h_\varphi(t_C)
 \le n|B_n|H_\varphi(a).
\]
Because $h_\varphi$ is continuous and has finite tail limsup,
$H_\varphi(a)<\infty$.  This proves finiteness and
\eqref{eq:envelope}.  Positivity follows because $K$ contains a centered ball
of positive radius and $\varphi$ is positive.

Let $S=\IS_\varphi(K)$ and choose $C_j\in\Cin(K)$ with
$\as_\varphi(C_j)\to S$.  By Blaschke's selection theorem
\cite{Schneider}, a subsequence converges in the Hausdorff metric to a compact
convex set $C_0\subset K$.  Suppose that $C_0$ has empty interior.  Then
$|C_j|\to0$, hence $t_j=\vrad(C_j)^{-2n}\to\infty$.  The preceding bound
gives
\begin{equation}\label{eq:collapse}
 \limsup_j\as_\varphi(C_j)
 \le n|B_n|\limsup_j h_\varphi(t_j)
 \le n|B_n|L_\varphi.
\end{equation}
Choose $\rho>0$ with $\rho B_n\subset K$.  By
$H_\varphi(\rho^{-2n})>L_\varphi$, there is a
$t\ge\rho^{-2n}$ with $h_\varphi(t)>L_\varphi$.  The ball
$rB_n$, $r=t^{-1/(2n)}\le\rho$, belongs to $\Cin(K)$, and therefore
\[
 S\ge\as_\varphi(rB_n)
 =n|B_n|h_\varphi(t)>n|B_n|L_\varphi,
\]
contradicting \eqref{eq:collapse}.  Thus $C_0$ has nonempty interior.
Centroid continuity gives $g(C_0)=0$, so $C_0\in\Cin(K)$.  Upper
semicontinuity now yields
\[
 S=\limsup_j\as_\varphi(C_j)
 \le\as_\varphi(C_0)\le S.
\]
Hence the supremum is attained.

For $K=RB_n$, every $t\ge R^{-2n}$ is realized by the centered ball
$t^{-1/(2n)}B_n\subset RB_n$.  Formula \eqref{eq:ball} supplies the reverse
inequality to \eqref{eq:envelope}, proving equality.  Notice also that the
supremum defining $H_\varphi(a)$ is attained at a finite $t$: otherwise a
maximizing sequence escaping to infinity would force
$H_\varphi(a)\le L_\varphi$, contrary to \eqref{eq:tailgap}.

It remains to prove continuity.  The functional $\IS_\varphi$ is increasing
under inclusion.  Also, for $c\ge1$, the source dilation estimate
$\as_\varphi(cC)\le c^n\as_\varphi(C)$ gives
\begin{equation}\label{eq:dilate}
 \IS_\varphi(cK)\le c^n\IS_\varphi(K).
\end{equation}
Suppose centered $K_j$ converges to centered $K$ in Hausdorff distance.
Since the origin is interior to $K$, standard support-function inclusions give
numbers $\delta_j,\eta_j\downarrow0$ such that, for all large $j$,
\[
 K_j\subset(1+\delta_j)K,
 \qquad K\subset(1+\eta_j)K_j.
\]
Monotonicity and \eqref{eq:dilate} imply
\[
 (1+\eta_j)^{-n}\IS_\varphi(K)
 \le\IS_\varphi(K_j)
 \le(1+\delta_j)^n\IS_\varphi(K),
\]
and continuity follows.

Finally, if $\varphi\in\operatorname{Conc}^{-}(0,\infty)$, then
$h_\varphi$ is positive and strictly decreasing.  Its finite tail limit is
$L_\varphi$, and
$H_\varphi(a)=h_\varphi(a)>L_\varphi$ for every $a>0$.  Thus the entire
source decreasing-ratio class is contained in the tail-gapped class.
\end{proof}

\section{A strict extension: the logarithmic function}

\begin{corollary}\label{cor:log}
Let $\varphi(t)=\log(1+t)$.  Then $\varphi$ is tail-gapped but belongs to
neither $\operatorname{Conc}^{-}(0,\infty)$ nor
$\operatorname{Conc}^{+}(0,\infty)$.  There is a unique $t_*>1$ satisfying
\begin{equation}\label{eq:tstar}
 \log(1+t_*)=\frac{2t_*}{1+t_*},
\end{equation}
and $t_*=3.9215536346\ldots$.  If
$h(t)=\log(1+t)/\sqrt t$, then
\[
 H(a)=
 \begin{cases}
 h(t_*),&0<a\le t_*,\\
 h(a),&a\ge t_*.
 \end{cases}
\]
Consequently Theorem \ref{thm:main} applies, and in particular
\[
 \IS_\varphi(B_n)=n|B_n|h(t_*)
 >n|B_n|h(1)=\as_\varphi(B_n).
\]
\end{corollary}

\begin{proof}
The function $\varphi$ is positive and concave, with $\varphi(0)=0$ and
$\varphi(t)/t\to0$.  Direct differentiation shows that the sign of $h'(t)$
is the sign of
\[
 q(t)=\frac{2t}{1+t}-\log(1+t).
\]
Here $q(0)=0$ and
\[
 q'(t)=\frac{1-t}{(1+t)^2}.
\]
Thus $q$ increases on $(0,1)$, decreases on $(1,\infty)$, and tends to
$-\infty$.  Since $q(1)=1-\log2>0$, it has a unique positive zero $t_*>1$.
Therefore $h$ increases up to $t_*$ and decreases afterwards.  In particular
it is globally nonmonotone.  Since $h(t)\to0$ as $t\to\infty$, it is
tail-gapped and the displayed envelope follows.  Equation \eqref{eq:tstar}
is $q(t_*)=0$; the reported decimal is checked reproducibly by
\texttt{code/check\_log\_example.py} and is not used in the proof.
\end{proof}

\begin{remark}[Repair of the source non-collapse step]
The proof of the source existence lemma writes that strict decrease of
$h_\varphi$ makes the right side of \eqref{eq:collapse} tend to zero.  This is
stronger than the definition.  For
$\varphi(t)=\sqrt t+t^{1/4}$, one has
$h_\varphi(t)=1+t^{-1/4}\downarrow1$.  The source conclusion is nevertheless
recovered by the strict tail gap: every finite $h_\varphi(t)$ lies strictly
above its tail limit.  Theorem \ref{thm:main} isolates exactly the compactness
mechanism needed here and also admits nonmonotone ratios.
\end{remark}

\section{Verification and limitations}

The proof was checked step by step against the following failure modes.
The affine isoperimetric bound is applied only to centered bodies; the tail
gap is imposed for every $a>0$, so it works for bodies of arbitrary inradius;
centroid continuity is invoked only after degeneration is excluded; and the
continuity argument uses only finite values of $\IS_\varphi$.  If
$L_\varphi=\infty$, arbitrarily small centered balls can make
$\IS_\varphi(K)=\infty$, so tail boundedness is a real restriction.  Without
the strict gap, this compactness proof need not prevent degeneration; no
necessity assertion is made.

The original problem is broad.  The theorem above extends existence,
attainment, continuity, and the radial affine-isoperimetric estimate for the
inner maximal functional.  The logarithmic example shows that the source's
ball identity and equality theory cannot extend unchanged.  A full extension
of the Blaschke-Santalo theorem remains open.

\paragraph{Bounded novelty check.}
The four local run indexes were searched by arXiv id, title, and core terms.
Exact-title and exact-phrase arXiv/web searches were run through 19 July 2026.
The journal page lists three citing works (a 2023 volume-product survey and two
2025 convex-geometry papers); none states the tail-envelope theorem.  The 2024
preprint \emph{Extremal affine surface areas in a functional setting}
(arXiv:2402.16130) treats a different functional construction and fixed power
parameters.  This bounded search found no prior statement of the present
extension, but absolute novelty is not certified.

\paragraph{Human review recommendation.}
Send to a convex-geometric analyst.  The highest-value checks are
\eqref{eq:collapse}, the exact hypotheses of the affine isoperimetric and
upper-semicontinuity inputs, and whether this theorem should be described as a
partial answer to the source's deliberately broad phrase ``the results.''

\begin{thebibliography}{9}
\bibitem{Hoehner}
S. Hoehner,
\emph{Extremal general affine surface areas},
J. Math. Anal. Appl. 505 (2022), 125506;
\href{https://arxiv.org/abs/2103.00294}{arXiv:2103.00294}.

\bibitem{Ludwig}
M. Ludwig,
\emph{General affine surface areas},
Adv. Math. 224 (2010), 2346--2360.

\bibitem{Schneider}
R. Schneider,
\emph{Convex Bodies: The Brunn--Minkowski Theory},
2nd ed., Cambridge University Press, 2014.
\end{thebibliography}

\end{document}
